\documentclass[11pt]{article}
\usepackage[letterpaper,margin=1in]{geometry}
\usepackage[T1]{fontenc}
\usepackage{lmodern,microtype}
\usepackage{amsmath,amssymb,amsthm,mathtools,bm}
\usepackage{booktabs,array,longtable}
\usepackage[dvipsnames]{xcolor}
\usepackage{hyperref}
\hypersetup{colorlinks=true,linkcolor=MidnightBlue,citecolor=MidnightBlue,urlcolor=MidnightBlue,
 pdftitle={RA-11: Partial results for unrestricted Kronecker trace estimation},
 pdfauthor={Sidney Holden},pdfsubject={Partial mathematical results, not a complete solution}}
\usepackage{fancyhdr}
\pagestyle{fancy}\fancyhf{}
\fancyhead[L]{\small RA-11: unrestricted Kronecker trace estimation}
\fancyhead[R]{\small Partial results}
\fancyfoot[C]{\thepage}
\setlength{\headheight}{14pt}
\setlength{\parskip}{4pt}
\setlength{\parindent}{0pt}
\allowdisplaybreaks
\newtheorem{theorem}{Theorem}[section]
\newtheorem{lemma}[theorem]{Lemma}
\newtheorem{corollary}[theorem]{Corollary}
\newtheorem{proposition}[theorem]{Proposition}
\theoremstyle{definition}\newtheorem{definition}[theorem]{Definition}
\theoremstyle{remark}\newtheorem{remark}[theorem]{Remark}
\newcommand{\R}{\mathbb R}\newcommand{\C}{\mathbb C}
\newcommand{\E}{\mathbb E}\newcommand{\PP}{\mathbb P}
\newcommand{\tr}{\operatorname{tr}}\newcommand{\rank}{\operatorname{rank}}
\newcommand{\Var}{\operatorname{Var}}\newcommand{\Cov}{\operatorname{Cov}}
\newcommand{\range}{\operatorname{range}}\newcommand{\diag}{\operatorname{diag}}
\newcommand{\poly}{\operatorname{poly}}
\newcommand{\eps}{\varepsilon}\newcommand{\Qprod}{Q_{\mathrm{prod}}}
\newcommand{\norm}[1]{\left\lVert#1\right\rVert}
\newcommand{\inner}[2]{\left\langle#1,#2\right\rangle}
\newcommand{\ceil}[1]{\left\lceil#1\right\rceil}
\newcommand{\floor}[1]{\left\lfloor#1\right\rfloor}
\newcommand{\contract}{\operatorname{contr}}
\begin{document}
\begin{center}
{\small RESEARCH NOTE \quad | \quad 13 SEPTEMBER 2026}\par\vspace{10pt}
{\LARGE\bfseries RA-11: unrestricted Kronecker\par trace estimation}\par\vspace{9pt}
{\large Polynomial real simulation, parallel factor access,\par and a bounded-probability lower bound}\par\vspace{10pt}
{\large Sidney Holden}\par\vspace{4pt}
{\small Center for Computational Biology, Flatiron Institute, Simons Foundation}\par\vspace{4pt}
{\small Prepared with ChatGPT; mathematical arguments and reproducible checks}
\end{center}
\vspace{8pt}
\noindent\fbox{\parbox{0.965\linewidth}{\textbf{Status: partial results, not a complete solution of RA-11.}
The general upper and lower bounds in Theorem~\ref{thm:main} do not match.
A tensor-product subclass is characterized up to constants, and the projection
conjecture identified below has an explicit counterexample. The arguments have
passed an independent informal Codex AI-agent audit of their stated partial
scopes; see the accompanying independent-review.md. This is not external human
peer review or formal theorem-prover verification.}}

\begin{abstract}
We study relative trace estimation of arbitrary real positive-semidefinite
matrices from adaptive, exact, real rank-one tensor matrix-vector queries.
A degree-$q$ interpolation identity simulates one complex product query with
$q+1$ real product queries. It also yields a counterexample, already for local
dimension two, to Conjecture~23 in arXiv:2502.08029v2: a fixed span of $q+1$
real product vectors has squared projection at least $2^{1-q}$ on every real
unit product vector.
For the promise $M=A_1\otimes\cdots\otimes A_q$, with PSD factors of order $n$,
a normalization-and-shift construction provides simultaneous ordinary
matrix-vector access to all factors. This gives exact recovery in $n$ original
queries and a relative-error estimator using $O(\sqrt q/\eps)$ queries.
An adaptive Wishart posterior argument, combined with log-concave
anti-concentration, gives a matching bounded-probability lower bound for this
subclass. Blocking yields a general lower bound, while interpolation and
ordinary-query simulation give general upper bounds. The remaining gap for
arbitrary PSD matrices is stated explicitly.
\end{abstract}
\newpage
\tableofcontents
\newpage

\section{Model, source audit, and main results}
\label{sec:main}

\subsection{The exact target}
Fix integers $n,q\ge2$, $0<\eps<1/2$, and $N=n^q$.
The unknown input is an arbitrary real symmetric PSD matrix
$M\in\R^{N\times N}$. An oracle call specifies $v_1,\ldots,v_q\in\R^n$
and returns the \emph{entire} vector
\[
 M(v_1\otimes\cdots\otimes v_q)\in\R^N.
\]
Let $Q(n,q,\eps)$ be the minimum worst-case number of calls needed by a
randomized adaptive algorithm whose output $T$ satisfies, for every fixed $M$,
\begin{equation}\label{eq:target}
 \PP\{|T-\tr M|\le\eps\tr M\}\ge 2/3.
\end{equation}
Arbitrary finite exact arithmetic between queries is allowed, and no condition
number bound is imposed. These are the quantifiers and oracle specified by
RA-11~\cite{ra11}. In particular, the matrix itself need not be a Kronecker
product. Throughout, real-valued Gaussian randomness is used in the usual
exact-real oracle model. All displayed query simulations require only finitely
many arithmetic operations on the sampled values and responses.

\subsection{What the cited sources do and do not supply}
Theorem~7 of Meyer--Swartworth--Woodruff~\cite{msw} concerns a conditioned
\emph{scalar quadratic-form} oracle. Its exponential lower bound is not a lower
bound for the stronger, unrestricted full-vector oracle above. Section~6 of
that paper poses Conjecture~23 about projection onto arbitrary spans of product
queries. Section~\ref{sec:counterexample} gives a counterexample to the exact
statement in arXiv v2; it does not contradict the paper's conditioned theorems.

Meyer--Avron~\cite{ma} analyze Kronecker-Hutchinson estimators, give a
root-mean-square-error lower bound using Wishart factors, and describe exact
recovery of tensor-product matrices in $qn+1$ calls. An estimator-specific
variance lower bound is not a minimax lower bound, and an RMSE lower bound alone
does not imply the success-probability criterion~\eqref{eq:target}. Our lower
bound below supplies a separate anti-concentration argument for that criterion.
The parallel access identity improves the stated $qn+1$ recovery construction
to $n$ calls under the PSD-factor promise. No claim of an exhaustive novelty
or priority check is made.

\subsection{General bounds and the gap}
Put
\begin{align}
 \alpha_n&=\frac{3n}{n+2}, & \beta_n&=\frac{2n}{n+1},\\
 L(n,q,\eps)&=\max_{1\le b\le q}
       \min\left\{n^b,\frac{\sqrt{\floor{q/b}}}{\eps}\right\}.
 \label{eq:L}
\end{align}
Define $U(n,q,\eps)$ to be the minimum of the following four integers:
\begin{align}
 U_0&=n^q,\label{eq:U0}\\
 U_{\R}&=\ceil{\frac{3(\alpha_n^q-1)}{\eps^2}},\label{eq:UR}\\
 U_{\C\to\R}&=(q+1)\ceil{\frac{3(\beta_n^q-1)}{\eps^2}},\label{eq:UC}\\
 U_{\rm ordinary}&=n^{q-1}\left(5\ceil{4/\eps}+4\right).
 \label{eq:Uordinary}
\end{align}

\begin{theorem}[General bounds]\label{thm:main}
For the full RA-11 model,
\begin{equation}\label{eq:mainbounds}
 \frac{1}{80000}L(n,q,\eps)
       \ \le\ Q(n,q,\eps)\ \le\ U(n,q,\eps).
\end{equation}
In particular, $Q(n,q,\eps)=\Theta(n^q)$ in the parameter regime
$0<\eps\le n^{-q}$, with universal constants.
\end{theorem}

The upper bound follows from Sections~\ref{sec:interpolation},
\ref{sec:upper}, and~\ref{sec:hutch}; the lower bound follows from
Sections~\ref{sec:lower} and~\ref{sec:blocking}. Taking $b=q$ in
\eqref{eq:L} proves the last assertion together with $U_0$.

\textbf{Unresolved in this note.} Away from such special regimes,
\eqref{eq:mainbounds} can have a very large gap. Neither a matching upper bound
for $L$ nor a matching lower bound for $U$ is established. Therefore
Theorem~\ref{thm:main} does not determine $Q$ up to universal constants
simultaneously in all three parameters.

\subsection{A subclass that is characterized}
For $n\ge2$ and $k\ge1$, let $\Qprod(n,k,\eps)$ denote the same minimax
quantity with the promise
\begin{equation}\label{eq:productpromise}
 M=A_1\otimes\cdots\otimes A_k,\qquad A_i\in\R^{n\times n},\quad A_i\succeq0.
\end{equation}
The local factors are unknown and are not supplied to the algorithm.

\begin{theorem}[Tensor-product subclass]\label{thm:productmain}
For all these parameters,
\begin{align}
 \frac1{80000}\min\left\{n,\frac{\sqrt k}{\eps}\right\}
 &\le\Qprod(n,k,\eps)\\
 &\le \min\left\{n,\;5\ceil{8\sqrt k/\eps}+5\right\}
 \le45\min\left\{n,\frac{\sqrt k}{\eps}\right\}.
 \label{eq:productmain}
\end{align}
The $n$-query upper bound recovers the entire matrix in factored form, and its
trace exactly, with probability one for every fixed promised input.
\end{theorem}

Only the subclass in~\eqref{eq:productpromise} has the matching rate above.
Low tensor-product complexity of the \emph{queries} does not imply this promise
about the \emph{matrix}.

\section{Complex product queries from \texorpdfstring{$q+1$}{q+1} real calls}
\label{sec:interpolation}

\begin{theorem}[Interpolation simulation]\label{thm:interpolation}
Let $M$ be any real linear map on $(\R^n)^{\otimes q}$, not necessarily symmetric
or PSD. For any $z_i=a_i+\mathrm{i}b_i\in\C^n$, the complexified response
$M(z_1\otimes\cdots\otimes z_q)$ is computable using $q+1$ real product
matrix-vector queries and finite exact arithmetic.
\end{theorem}
\begin{proof}
Consider the vector-valued polynomial
\[
 p(t)=\bigotimes_{i=1}^q(a_i+t b_i),\qquad \deg p\le q.
\]
Choose distinct real numbers $t_0,\ldots,t_q$ and set
\[
 \ell_j(t)=\prod_{\substack{0\le h\le q\\h\ne j}}
                  \frac{t-t_h}{t_j-t_h}.
\]
Coordinatewise Lagrange interpolation, followed by real linearity, gives
\begin{equation}\label{eq:interpolate}
 M\left(\bigotimes_{i=1}^q z_i\right)
    =M p(\mathrm{i})
    =\sum_{j=0}^q\ell_j(\mathrm{i})\,
             M\left(\bigotimes_{i=1}^q(a_i+t_jb_i)\right).
\end{equation}
Every actual oracle input has real factors. Complex numbers on the right are
stored as pairs of reals. Integer nodes make all interpolation coefficients
Gaussian-rational, so even the coefficients require no limiting operation or
root computation.
\end{proof}

The simulation is adaptive: after reconstructing a complex response exactly,
one may form the next complex query of an arbitrary adaptive algorithm and
repeat. Thus a $t$-query complex-product algorithm on a real matrix gives a
$(q+1)t$-query real-product algorithm with the same output distribution.

\begin{remark}[Conditioning and arithmetic scope]
The interpolation coefficients may be large, and the real query matrix can be
very ill-conditioned. The proof is an exact identity, not a stable numerical
algorithm. It neither uses infinitesimal inputs nor extracts digits of unknown
real numbers. An expansion with $2^q$ real product terms is always possible but
is not necessary in this unrestricted exact model.
\end{remark}

\section{An explicit counterexample to the projection conjecture}
\label{sec:counterexample}

Conjecture~23 of arXiv:2502.08029v2~\cite{msw} asserts an exponentially small
projection bound for a random real unit product vector onto a fixed span of
polynomially many real product vectors. The particular asserted scale is
$c_1^{-q}/n^q$, with probability at least $1-c_2^{-q}$, for constants
$c_1,c_2>1$. The following deterministic inequality contradicts that scale.

\begin{theorem}[Failure at $n=2$]\label{thm:counterexample}
For every $q\ge2$ there is a fixed $(q+1)$-dimensional subspace
$S\subseteq(\R^2)^{\otimes q}$ spanned by $q+1$ real product vectors such that
\begin{equation}\label{eq:projectionlower}
 \norm{P_Su}_2^2\ge 2^{1-q}
\end{equation}
for every real unit product vector $u=u_1\otimes\cdots\otimes u_q$.
\end{theorem}
\begin{proof}
Choose distinct real nodes $t_0,\ldots,t_q$ and let
\[
 S=\operatorname{span}_{\R}\{(e_1+t_j e_2)^{\otimes q}:0\le j\le q\}.
\]
The coefficient vectors $D_h$ in
\[
 (e_1+t e_2)^{\otimes q}=\sum_{h=0}^q t^h D_h
\]
have disjoint nonempty supports: $D_h$ is the sum of coordinate tensors with
exactly $h$ occurrences of $e_2$. They are linearly independent. The Vandermonde
matrix at the distinct nodes is nonsingular, hence $\dim S=q+1$.

Set
\[
 z=(e_1+\mathrm{i}e_2)^{\otimes q},\qquad r=\operatorname{Re}z,
       \qquad s=\operatorname{Im}z.
\]
Interpolation places both $r$ and $s$ in $S$. Their supports have respectively
even and odd parity, so
\[
 r^Ts=0,\qquad \norm r_2^2=\norm s_2^2=2^{q-1}.
\]
Write $u_i=(a_i,b_i)^T$, with $a_i^2+b_i^2=1$. Then
\[
 (r^Tu)^2+(s^Tu)^2
   =\left|\prod_{i=1}^q(a_i+\mathrm{i}b_i)\right|^2
   =\prod_{i=1}^q(a_i^2+b_i^2)=1.
\]
Projection onto the two-dimensional subspace spanned by $r,s$ therefore has
squared norm exactly $2^{1-q}$. Projection onto its superspace $S$ can only
increase that norm.
\end{proof}

For any $c_1>1$,
\[
 2^{1-q}>c_1^{-q}2^{-q}.
\]
Thus the conjectured upper-bound event has probability zero for this fixed
choice of $S$, rather than probability at least $1-c_2^{-q}>0$.
The number of queries is $q+1$, which is polynomial; normalizing each spanning
product vector leaves $S$ unchanged. This contradicts the universal polynomial-
size-span formulation in the cited v2. It does not exclude restricted statements
with an explicit condition-number bound or a substantially smaller allowed span.


\begin{proposition}[An explicit conditioning calculation]\label{prop:conditioning}
The same subspace $S$ has a normalized real product spanning family $V$ with
\[
 \kappa_2(V)^2=\binom q{\floor{q/2}}\ge\frac{2^q}{q+1}.
\]
Thus even unit-norm query columns can give the exponentially ill-conditioned
span used by the counterexample.
\end{proposition}
\begin{proof}
Let $L=q+1$, $\theta_j=\pi j/L$, and
$v_j=(\cos\theta_j\,e_1+\sin\theta_j\,e_2)^{\otimes q}$ for $0\le j<L$.
The Gram matrix has entries
\[
 (V^TV)_{jh}=\cos^q(\theta_j-\theta_h)
 =2^{-q}\sum_{a=0}^q\binom qa
          e^{\mathrm{i}(q-2a)(\theta_j-\theta_h)}.
\]
The matrix $F_{ja}=L^{-1/2}e^{\mathrm{i}(q-2a)\theta_j}$ is unitary, by the
finite geometric-series identity. Hence the Gram eigenvalues are exactly
$L2^{-q}\binom qa$, $0\le a\le q$. They are all positive, so the columns span
the full symmetric tensor subspace $S$, and their largest-to-smallest ratio is
the stated binomial coefficient. Its lower bound follows because the largest
of the $q+1$ binomial coefficients is at least their average $2^q/(q+1)$.
\end{proof}

\section{General-purpose trace-estimation upper bounds}
\label{sec:upper}

\subsection{Real sphere product probes}
Let $x_i$ be independent uniform vectors on the real sphere of radius $\sqrt n$,
and let $x=\bigotimes_i x_i$. Then $\E[xx^T]=I_N$.
For a local coefficient vector $a$,
\[
 \E(a^Tx_i)^4=\frac{3n}{n+2}\norm a_2^4=\alpha_n\norm a_2^4.
\]
We record an induction that is useful because it requires no assumptions on the
eigenvectors of $M$.

\begin{lemma}\label{lem:real4}
For every $a\in(\R^n)^{\otimes q}$,
\[
 \E|a^Tx|^4\le\alpha_n^q\norm a_2^4.
\]
\end{lemma}
\begin{proof}
The local identity proves the case $q=1$. For the induction step, write the
contraction with the last $q-1$ factors as $f=(f_1,\ldots,f_n)^T$. Conditioning
on those factors gives $\E_{x_1}(f^Tx_1)^4=\alpha_n(\sum_j f_j^2)^2$.
If $a_j$ is the corresponding coefficient tensor for $f_j$, Cauchy--Schwarz
and the induction hypothesis give
\[
 \E\Big(\sum_j f_j^2\Big)^2
 \le\sum_{j,h}(\E f_j^4\,\E f_h^4)^{1/2}
 \le\alpha_n^{q-1}\Big(\sum_j\norm{a_j}_2^2\Big)^2.
\]
Multiplication by the first factor $\alpha_n$ proves the assertion.
\end{proof}

Let $M=\sum_j\lambda_j a_ja_j^T$ be a spectral decomposition with
$\lambda_j\ge0$ and unit $a_j$, and put $\tau=\tr M$.
Cauchy--Schwarz in $L^2$, or the triangle inequality in $L^2$, gives
\[
 \E(x^TMx)^2
 \le\left(\sum_j\lambda_j\sqrt{\E(a_j^Tx)^4}\right)^2
 \le\alpha_n^q\tau^2.
\]
Consequently the average of $m$ independent scalar observations has variance
at most $(\alpha_n^q-1)\tau^2/m$. Chebyshev's inequality with
$m=U_{\R}$ proves~\eqref{eq:UR}. Each scalar is computed from one full-vector
oracle response. For $M=0$ the output is exactly zero.

\subsection{Complex sphere product probes, implemented with real inputs}
Let $z_i$ be independent uniform complex sphere vectors of radius $\sqrt n$,
let $z=\bigotimes_i z_i$, and put $Z=z^*Mz$. The usual partial trace
$\tr_S M$ traces out subsystems indexed by $S$; for $S=[q]$ its value is the
scalar $\tau$, viewed as a $1\times1$ matrix.

\begin{lemma}[Exact second moment]\label{lem:complexmoment}
For every Hermitian PSD $M$,
\begin{align}
 \E Z&=\tau,\\
 \E Z^2&=\left(\frac n{n+1}\right)^q
       \sum_{S\subseteq[q]}\norm{\tr_S M}_F^2
       \le\beta_n^q\tau^2.
 \label{eq:complexmoment}
\end{align}
The upper bound is attained by a rank-one projector onto a product vector.
\end{lemma}
\begin{proof}
The local second and fourth moments are
\[
 \E[z_i z_i^*]=I_n,\qquad
 \E[(z_i z_i^*)\otimes(z_i z_i^*)]
       =\frac n{n+1}(I+F_i),
\]
where $F_i$ swaps the two copies of local subsystem $i$.
Expanding $\prod_i(I+F_i)$ and contracting with $M\otimes M$ gives
\eqref{eq:complexmoment}; the swap indexed by $S$ gives
$\tr[(\tr_{S^c}M)^2]$, and complementation only permutes the sum.
Every partial trace is PSD and has trace $\tau$, so its squared Frobenius norm
is at most $\tau^2$. There are $2^q$ terms. For a product rank-one matrix every
partial trace is rank one or a scalar, so all these inequalities are equalities.
\end{proof}

This is the complex sphere version of the partial-trace moment method
in~\cite{ma}. Averaging
$m=\ceil{3(\beta_n^q-1)/\eps^2}$ independent copies and applying Chebyshev gives
failure probability at most $1/3$. Theorem~\ref{thm:interpolation} computes each
complex response in $q+1$ real calls, proving~\eqref{eq:UC} for the real oracle.

No square root is needed to implement sphere normalization. Draw
$w_i=a_i+\mathrm{i}b_i$ with independent real standard Gaussian entries,
simulate $M(\bigotimes_i w_i)$, and return
\[
 \left(\prod_{i=1}^q\frac n{\norm{w_i}_2^2}\right)
       (\bigotimes_i w_i)^*M(\bigotimes_i w_i).
\]
This has precisely the same distribution as $Z$, and the zero-denominator
exceptions have probability zero.

\subsection{Coordinate and ordinary-query simulations}
Querying every coordinate product vector reveals every diagonal entry and gives
$U_0=N$ calls. For an arbitrary vector $g\in\R^{n^q}$, reshape it as a tensor and
write
\begin{equation}\label{eq:fibredecomp}
 g=\sum_{j_2,\ldots,j_q=1}^n
      g_{:,j_2,\ldots,j_q}\otimes e_{j_2}\otimes\cdots\otimes e_{j_q}.
\end{equation}
Thus any ordinary matrix-vector product $Mg$ costs at most $n^{q-1}$ product
queries. Lemma~\ref{lem:hutch} below, with $r=\ceil{4/\eps}$, supplies an ordinary
algorithm with $5r+4$ calls and failure probability at most $1/4$.
Combining it with~\eqref{eq:fibredecomp} proves~\eqref{eq:Uordinary}.

\section{Exact recovery under a low-rank PSD promise}
\label{sec:lowrank}

\begin{proposition}\label{prop:lowrank}
If $M\succeq0$ has known rank bound $r$, then $r$ independent generic real
product queries recover $M$ exactly with probability one.
\end{proposition}
\begin{proof}
Let $s=\rank M\le r$. The images under $M$ of the coordinate product basis span
$\range M$, so some $s$ such images are independent. Therefore the determinant
of an appropriate $s\times s$ minor of $[Mx_1,\ldots,Mx_r]$, considered as a
polynomial in all the local factor entries of $x_1,\ldots,x_r$, is not the zero
polynomial. Independent Gaussian entries avoid its zero set almost surely.
Consequently $Y=MX$, where $X=[x_1,\ldots,x_r]$, has rank $s$ almost surely.

Select $s$ independent columns $Y_J$, with corresponding input columns $X_J$.
The matrix $K=X_J^TY_J$ is positive definite. To verify this and the reconstruction,
write $M=CC^T$ with $C\in\R^{N\times s}$ full column rank and set $D=C^TX_J$.
Then $Y_J=CD$ and $D$ is invertible. Hence
\[
 K=D^TD\succ0,\qquad
 Y_JK^{-1}Y_J^T=CD(D^TD)^{-1}D^TC^T=CC^T=M.
\]
For $s=0$ all responses vanish and the reconstruction is the zero matrix.
Rank selection and matrix inversion can be performed by finite exact arithmetic.
\end{proof}

In particular, for a nonzero rank-one PSD matrix and a generic query $x$ with
response $y=Mx$,
\begin{equation}\label{eq:rankone}
 \tr M=\frac{\norm y_2^2}{x^Ty}.
\end{equation}
This one-query fact is also noted in~\cite{ma}. It illustrates why a scalar
quadratic-form lower bound cannot simply be relabeled a full-vector lower bound.

\section{Parallel access to all PSD tensor factors}
\label{sec:parallel}

Assume the product promise~\eqref{eq:productpromise}. The following construction
is an exact oracle reduction; it does not require identifying local factors by
a numerical tensor decomposition.

Choose independent standard Gaussian reference vectors $g_i\in\R^n$ and query
\[
 y_0=M(g_1\otimes\cdots\otimes g_k).
\]
Compute the scalar
\[
 \gamma=(g_1\otimes\cdots\otimes g_k)^Ty_0.
\]
For the proof only, define
\begin{equation}\label{eq:gamma}
 \gamma_i=g_i^TA_i g_i,\qquad B_i=A_i/\gamma_i.
\end{equation}
If $M\ne0$, every $A_i$ is nonzero and PSD, so all $\gamma_i>0$ almost surely.
Then
\[
 \gamma=\prod_i\gamma_i>0,\qquad M=\gamma\bigotimes_i B_i,
       \qquad g_i^TB_i g_i=1.
\]
The individual $\gamma_i$ need not be known. If $\gamma=0$, return zero; this is
correct on $M=0$ and the exceptional event has probability zero on every fixed
nonzero promised input.

For a tensor $y$, let $\contract_{-i}(y;g)$ mean contraction of every output mode
$j\ne i$ with $g_j$. Compute
\begin{equation}\label{eq:bi}
 b_i=\frac{\contract_{-i}(y_0;g)}{\gamma}=B_i g_i.
\end{equation}
The equality follows because each omitted contraction contributes $\gamma_j$.
In particular, $b_i^Tg_i=1$.

\begin{theorem}[One round of all local products]\label{thm:parallel}
After the reference query, any prescribed tuple of vectors
$v_1,\ldots,v_k\in\R^n$ can be mapped to the entire tuple
$(B_1v_1,\ldots,B_kv_k)$ with one further original product query.
The vectors can depend on all previous responses.
\end{theorem}
\begin{proof}
Using the known $b_i$, define
\begin{equation}\label{eq:shift}
 \alpha_i=1-b_i^Tv_i,\qquad w_i=v_i+\alpha_i g_i.
\end{equation}
Because $B_i$ is symmetric and $b_i=B_i g_i$,
\[
 g_i^TB_iw_i=b_i^Tw_i=b_i^Tv_i+\alpha_i b_i^Tg_i=1.
\]
Query $y=M(\bigotimes_i w_i)$. All other-mode contractions are now normalized,
so
\[
 \frac{\contract_{-i}(y;g)}\gamma=B_iw_i.
\]
Subtracting $\alpha_i b_i$ recovers $B_i v_i$. No division by an unknown
bilinear form is used, even when $b_i^Tv_i=0$ or $B_iv_i=0$.
\end{proof}

\begin{corollary}[Exact $n$-call recovery]\label{cor:exactn}
Every promised matrix~\eqref{eq:productpromise} can be recovered, in the form
$\gamma\bigotimes_i B_i$, using at most $n$ original calls almost surely.
\end{corollary}
\begin{proof}
One local column $B_i g_i=b_i$ is already known. Complete $g_i$ to a basis using
$n-1$ coordinate vectors. Such a completion can be found by omitting any
coordinate $e_j$ for which $(g_i)_j\ne0$.
Use Theorem~\ref{thm:parallel} for the $n-1$ remaining basis columns in parallel
across all $i$. If $X_i$ is the resulting known local basis and $Y_i=B_iX_i$ its
known image, $B_i=Y_iX_i^{-1}$. The total is one reference call plus $n-1$ rounds.
Finally, $\tr M=\gamma\prod_i\tr B_i$.
\end{proof}

The promise is essential: for a general PSD matrix, contractions of a response
do not isolate independent local operators, and the normalization in
Theorem~\ref{thm:parallel} is not an oracle reduction for that general matrix.

\section{A local unbiased Hutch++ lemma and the product upper bound}
\label{sec:hutch}

The following version of the standard Hutch++ construction~\cite{hutchpp} is
included to make the constants, unbiasedness, and scale-equivariance used in the
parallel product estimator explicit.

\begin{lemma}[Ordinary PSD trace estimation]\label{lem:hutch}
Let $A\in\R^{d\times d}$ be PSD and $r\ge1$ an integer. There is a nonnegative,
unbiased estimator $H_r(A)$ using at most $5r+4$ ordinary matrix-vector calls,
such that
\begin{equation}\label{eq:hutchvar}
 \Var H_r(A)\le\frac4{r(r+1)}(\tr A)^2
              \le\frac4{r^2}(\tr A)^2.
\end{equation}
It can be chosen scale-equivariant for every fixed realization of its randomness:
$H_r(cA)=cH_r(A)$ for every $c>0$.
\end{lemma}
\begin{proof}
If $r\ge d$, query a basis and return the exact trace. Otherwise let
$m=2r+2$, draw a standard Gaussian $d\times m$ matrix $\Omega$, query
$Y=A\Omega$, and let $P$ be the orthogonal projection onto $\range Y$.
These are $m$ calls. Using at most $m$ further calls, compute $\tr(PA)$ exactly.
For example, choose independent columns $F$ of $Y$ and use
\[
 P=F(F^TF)^{-1}F^T,\qquad
 \tr(PA)=\tr\big((F^TF)^{-1}F^TAF\big).
\]
This formula avoids a square-root operation in the implementation.

Let $R=(I-P)A(I-P)\succeq0$ and draw $r$ fresh independent standard Gaussian
vectors $h_\ell$. With one call to $A(I-P)h_\ell$ for each, form
\[
 H_r(A)=\tr(PA)+\frac1r\sum_{\ell=1}^r h_\ell^T R h_\ell.
\]
Conditioned on $P$, this is nonnegative, has mean $\tr A$, and has variance
$2\norm R_F^2/r$. Its total call count is $2m+r=5r+4$.

It remains to bound the expected residual norm. Diagonalize $A$, partition its
singular values after the leading $r$, and split the Gaussian matrix as
$\Omega=(\Omega_1^T,\Omega_2^T)^T$, where $\Omega_1$ has $r$ rows.
It has full row rank almost surely. The matrix with columns
$Y\Omega_1^\dagger[I_r\ \ 0]$ lies in $\range Y$ and agrees with the leading
block of $A$. Its lower-block error shows
\[
 \norm{(I-P)A}_F^2
 \le \sum_{j>r}\lambda_j^2
        +\norm{\Lambda_2\Omega_2\Omega_1^\dagger}_F^2.
\]
Conditioning on $\Omega_1$ and averaging $\Omega_2$ makes the second term equal
in expectation to $(\sum_{j>r}\lambda_j^2)\norm{\Omega_1^\dagger}_F^2$.
Moreover,
\begin{equation}\label{eq:inversewishart}
 \E\norm{\Omega_1^\dagger}_F^2=\frac r{m-r-1}<1.
\end{equation}
For completeness, a diagonal entry of $(\Omega_1\Omega_1^T)^{-1}$ is the
reciprocal squared norm of one row projected orthogonally to the other $r-1$
rows. Conditional on those rows the squared norm has distribution
$\chi^2_{m-r+1}$. Integrating its gamma density gives reciprocal mean
$1/(m-r-1)$; summing the diagonal proves~\eqref{eq:inversewishart}.

Since $\lambda_{r+1}\le\tr A/(r+1)$,
\[
 \sum_{j>r}\lambda_j^2
 \le\lambda_{r+1}\sum_{j>r}\lambda_j
 \le\frac{(\tr A)^2}{r+1}.
\]
Also $\norm R_F\le\norm{(I-P)A}_F$. Therefore
\[
 \E\norm R_F^2\le\frac{2(\tr A)^2}{r+1}.
\]
Taking the expectation of the conditional variance gives~\eqref{eq:hutchvar};
the conditional mean is constant, so there is no additional variance term.
Scaling $A$ by $c>0$ preserves $P$ for fixed $\Omega$ and scales both estimator
terms by $c$, proving scale-equivariance.
\end{proof}

For ordinary relative error $\eps$, choosing $r=\ceil{4/\eps}$ gives failure
probability at most $1/4$ by Chebyshev, as used earlier.

\begin{proposition}[Parallel product estimator]\label{prop:productupper}
Under~\eqref{eq:productpromise}, at most
$5\ceil{8\sqrt k/\eps}+5$ original queries suffice for
relative error $\eps$ with probability at least $7/8$.
\end{proposition}
\begin{proof}
Perform the reference normalization of Section~\ref{sec:parallel} and use
Theorem~\ref{thm:parallel} to run the local algorithms $H_r(B_i)$ simultaneously,
with independent randomness for each factor and $r=\ceil{8\sqrt k/\eps}$.
They require at most $5r+4$ rounds after the reference call. Output
\[
 T=\gamma\prod_{i=1}^k H_r(B_i).
\]
Write $\tau_i=\tr A_i$ and first suppose $M\ne0$.
Scale-equivariance gives the pathwise identity
\[
 \frac{H_r(B_i)}{\tr B_i}
   =\frac{H_r(A_i)}{\tr A_i}.
\]
Thus the normalized local outputs depend only on their independent local
algorithmic randomness, not on the shared reference normalization. They are
independent, unbiased for one, and have variances at most $4/r^2$.
Since $\tr M=\gamma\prod_i\tr B_i$,
\[
 \E T=\tr M,\qquad
 \frac{\Var T}{(\tr M)^2}
 \le\left(1+\frac4{r^2}\right)^k-1
 \le e^{4k/r^2}-1
 \le\frac{\eps^2}{8}.
\]
The last inequality uses $4k/r^2\le\eps^2/16<1/64$ and
$e^x-1\le2x$ for $0\le x\le1$. Chebyshev gives failure at most $1/8$.
The zero input is handled by the reference query.
\end{proof}

Taking the smaller of this algorithm and Corollary~\ref{cor:exactn} proves the
upper bounds in Theorem~\ref{thm:productmain}.

\section{An adaptive, bounded-probability lower bound}
\label{sec:lower}

\subsection{A stronger parallel oracle}
For a lower bound on the product subclass, grant the algorithm a stronger oracle:
one round returns the individual vectors $A_i v_i$ for all $i$, rather than only
their tensor product. Permit arbitrary local vectors and adaptation across
factors, even within a round. Equivalently, allow at most $t$ ordinary queries
per factor, in any adaptive interleaving. A product-oracle algorithm with $t$
calls is simulated in this stronger model by taking the tensor product of the
individual responses. A lower bound in the stronger model therefore transfers
to the original one.

The hard prior uses independent square Wishart matrices in a known local
subspace. No condition-number assumption is imposed on the algorithm's queries.

\subsection{Adaptive Wishart posterior}

\begin{lemma}[Residual trace distribution]\label{lem:posterior}
Let $W_i=G_i^TG_i$, where the $G_i\in\R^{d\times d}$ have independent standard
Gaussian entries, independently over $1\le i\le k$.
After any adaptively interleaved transcript with at most $t<d$ ordinary
matrix-vector queries to each factor, one can augment the transcript with
additional queries so that each factor has $t$ independent queried directions.
Conditional on this augmented transcript and all algorithmic randomness,
\begin{equation}\label{eq:posterior}
 \tr W_i=a_i+X_i,\qquad X_i\stackrel{\rm independent}{\sim}\chi^2_{(d-t)^2},
\end{equation}
where $a_i$ is measurable from the transcript and satisfies
$0\le a_i\le\tr W_i$ pathwise. The conditional laws of the $X_i$ do not depend
on the transcript.
\end{lemma}
\begin{proof}
We give the Gaussian block calculation and then the adaptive induction. For a
fixed $t$-dimensional queried coordinate subspace, write $G=[G_1\ G_2]$, with
$G_1$ having $t$ columns. The observations give
\[
 K=G_1^TG_1\succ0,\qquad C=G_1^TG_2.
\]
Conditional on $G_1$, split $G_2$ into its orthogonal projections onto
$\range G_1$ and its orthogonal complement. Gaussian orthogonal projections
are independent. The first part is $G_1K^{-1}C$; in an orthonormal basis for the
complement, the second part is a $(d-t)\times(d-t)$ standard Gaussian matrix
$H$. Consequently
\begin{equation}\label{eq:schur}
 W=\begin{bmatrix}K&C\\C^T&C^TK^{-1}C+H^TH\end{bmatrix},
\end{equation}
with $H^TH$ independent of $(K,C)$. This independence also holds after integrating
out $G_1$, because its conditional law is the same for every $G_1$.
The trace is the sum of
\[
 a=\tr K+\tr(C^TK^{-1}C)
\]
and $\tr(H^TH)\sim\chi^2_{(d-t)^2}$. Both summands are nonnegative.

To justify adaptivity, maintain the following invariant after $s$ independent
queries to a factor. In a transcript-measurable orthonormal basis of its queried
span and complement, the still unknown Schur complement has law
$H_s^TH_s$, for a $(d-s)\times(d-s)$ standard Gaussian $H_s$, independently of
all current observations and of the other factors' unknown Schur complements.
The invariant holds initially. Subtract the known queried-span component from
the next vector. If the residual is zero, its response was already determined
and reveals no new information. Otherwise its normalized direction is a
measurable function of the current transcript. Conditional on that transcript,
rotate the complement to make this direction the first basis vector. Orthogonal
invariance leaves the Schur-complement distribution unchanged. The newly
observed column is precisely one column of that Schur complement plus known
terms. Applying~\eqref{eq:schur} with one queried column removes a row and a
column, and leaves a fresh square Wishart Schur complement of order $d-s-1$,
independent of the new observations. Conditional independence for other factors
is preserved, since the new response depends only on this factor's residual
and the already fixed transcript. This proves the invariant for arbitrary
interleavings by induction.

Dependent queries can be answered from prior responses. At the end, query fresh
orthogonal directions to pad each local span to dimension $t$; this only gives
the algorithm additional information. The invariant now proves
\eqref{eq:posterior}. Coordinate-free, the known nonnegative term is
\[
 a_i=\tr\left(W_i U_i(U_i^TW_iU_i)^{-1}U_i^TW_i\right),
\]
where $U_i$ has orthonormal columns spanning the final queried subspace.
It is observable from $W_iU_i$, and the remaining Schur-complement trace is
nonnegative. This proves $0\le a_i\le\tr W_i$.
\end{proof}

The augmentation is a device for bounding success from above: an estimator
based on the original transcript is still a measurable estimator given the
augmented transcript. It is not assumed that the original algorithm actually
used all the extra responses.

\subsection{Three anti-concentration facts}

\begin{lemma}[Logarithms of shifted chi-squares]\label{lem:loggamma}
Let $X\sim\chi^2_\nu$, with $\nu\ge2$, and $a\ge0$. Then $Y=\log(a+X)$ has a
log-concave density and
\begin{equation}\label{eq:logvar}
 \Var Y\ge\frac{2\nu}{(a+\nu+2)^2}.
\end{equation}
\end{lemma}
\begin{proof}
On its support, its log-density, up to a constant, is
\[
 y+\left(\frac\nu2-1\right)\log(e^y-a)-\frac{e^y-a}{2}.
\]
The second derivative equals
\[
 -\left(\frac\nu2-1\right)\frac{a e^y}{(e^y-a)^2}-\frac{e^y}{2}\le0.
\]
The support is an interval; extending the density by zero preserves
log-concavity.
For the variance, integrating the chi-square density by parts gives
\[
 \Cov(X,\log(a+X))=2\E\frac X{a+X}.
\]
The boundary terms vanish; the case $a=0$ follows either directly or by a limit.
Cauchy--Schwarz applied to
$X=\sqrt{X/(a+X)}\sqrt{X(a+X)}$ yields
\[
 \E\frac X{a+X}\ge
 \frac{(\E X)^2}{\E[X(a+X)]}
 =\frac\nu{a+\nu+2}.
\]
Since $\Var X=2\nu$, a second Cauchy--Schwarz inequality, this time for the
covariance, proves~\eqref{eq:logvar}.
\end{proof}

\begin{lemma}[One-dimensional log-concave density bound]\label{lem:lc}
A log-concave probability density $f$ on the real line with variance
$\sigma^2>0$ satisfies
\[
 \norm f_\infty\le\frac8\sigma.
\]
Thus the probability of any interval of length $\delta$ is at most
$8\delta/\sigma$.
\end{lemma}
\begin{proof}
Let $H=\norm f_\infty$ and let $m$ be a mode; limiting arguments handle an
endpoint mode. Let $a,b\ge0$ be the distances from $m$ to the two endpoints of
the superlevel interval where $f\ge H/e$. Its mass implies
$a+b\le e/H$. Log-concavity implies the tail bound
$f(m+t)\le H e^{-t/b}$ for $t\ge b$, and the corresponding left-tail bound
with $a$. A side with zero width contributes zero; a truncated support only
reduces the integral. Therefore
\[
 \int (x-m)^2 f(x)\,dx
 \le H(a^3+b^3)\left(\frac13+\frac5e\right)
 \le\frac{e^3/3+5e^2}{H^2}<\frac{64}{H^2}.
\]
Here $\int_1^\infty u^2e^{-u}\,du=5/e$. The variance is at most the second
moment about $m$, so $H\sigma\le8$.
\end{proof}

We also use the standard fact that a convolution of integrable log-concave
functions is log-concave, hence a sum of independent random variables with
log-concave densities has a log-concave density. This is a consequence of
Pr\'ekopa's marginal theorem; see the review~\cite{logconcave}. Its hypotheses
hold for all the densities in Lemma~\ref{lem:loggamma}.

\subsection{Success probability under the hard prior}

\begin{proposition}[Parallel oracle lower bound]\label{prop:lower}
For $n\ge2$, $k\ge1$, and $0<\eps<1/2$, every randomized algorithm succeeding
with probability at least $2/3$ for every tuple of $n\times n$ PSD matrices, when
estimating $\prod_i\tr A_i$ from the stronger parallel oracle, needs at least
\[
 \frac1{80000}\min\left\{n,\frac{\sqrt k}{\eps}\right\}
\]
queries per factor in the worst case.
\end{proposition}
\begin{proof}
We first work in dimension $d\ge4$ and suppose the local query budget is
$t\le d/4$. Fix the algorithm's random seed and use the independent Wishart
prior of Lemma~\ref{lem:posterior}; alternatively include that seed in the
conditioning throughout. After augmentation, write
\[
 \tr W_i=a_i+X_i,\qquad X_i\sim\chi^2_\nu,\qquad\nu=(d-t)^2.
\]
The known $a_i$ are nonnegative and $\E\sum_i a_i\le kd^2$, because
$\E\tr W_i=d^2$. Thus the transcript event
\[
 \mathcal G=\left\{\sum_i a_i\le10kd^2\right\}
\]
has probability at least $9/10$ by Markov. On $\mathcal G$, at least $k/2$ of
the indices satisfy $a_i\le20d^2$. Also
$9d^2/16\le\nu\le d^2$. For every such index,
Lemma~\ref{lem:loggamma} gives
\[
 \Var(\log(a_i+X_i)\mid\text{transcript})
 \ge\frac{2\nu}{(a_i+\nu+2)^2}\ge\frac1{500d^2}.
\]
For the last inequality it suffices to use
$a_i+\nu+2\le22d^2$ and $2\nu\ge9d^2/8$.

Conditioned on the augmented transcript, the log of the desired trace product is
\[
 Z=\log\prod_i\tr W_i=\sum_i\log(a_i+X_i).
\]
It has a log-concave density, and on $\mathcal G$ its conditional variance is at
least $k/(1000d^2)$. Lemma~\ref{lem:lc} therefore bounds its conditional density
by
\[
 8\sqrt{1000}\,\frac d{\sqrt k}.
\]
For a fixed positive estimate $\widehat T$, relative-error success requires
\[
 Z\in\left[\log\frac{\widehat T}{1+\eps},\,
             \log\frac{\widehat T}{1-\eps}\right].
\]
The interval has length
$\log((1+\eps)/(1-\eps))\le3\eps$ for $0<\eps<1/2$.
A nonpositive estimate has zero success probability under this strictly positive
prior. The same bound holds for randomized output after conditioning on all
its randomness. Consequently every estimator satisfies
\begin{equation}\label{eq:successbound}
 \PP\{\text{relative-error success}\}
 \le\frac1{10}+24\sqrt{1000}\,\frac{\eps d}{\sqrt k}.
\end{equation}
This is a probability bound, not merely a lower bound on mean-square error.

To obtain a uniform bound in the original local dimension $n$, put
$X=\sqrt k/\eps$ and choose
\[
 d=\min\{n,\floor{X/10000}\}.
\]
When $d\ge4$, embed each $d\times d$ Wishart matrix in the first $d$ coordinates
of an $n\times n$ matrix and fill the rest with zeros. The algorithm can be told
this support; queries outside it reveal nothing else. For $t\le d/4$, the
right-hand side of~\eqref{eq:successbound} is at most
\[
 \frac1{10}+\frac{24\sqrt{1000}}{10000}<0.176<\frac23.
\]
An algorithm that succeeds on every fixed input would have average success at
least $2/3$ under this prior, a contradiction. Thus $t>d/4$ is necessary.
Since $d\ge\min(n,X)/20000$ when $d\ge4$, this gives the desired constant.

If $d<4$, then either $n<4$ or $X<40000$, so $\min(n,X)<40000$.
Zero calls cannot succeed for every trace: the required output events for a
zero input and for a fixed positive-trace input are disjoint and would each
need probability at least $2/3$. Hence at least one call is necessary, which is
more than $\min(n,X)/80000$ in this remaining case.
The prior and every estimate above are independent of the random seed chosen
by the algorithm, so averaging over that seed preserves the argument.
\end{proof}

The stronger parallel oracle simulates the original product oracle, so
Proposition~\ref{prop:lower} supplies the lower bound of
Theorem~\ref{thm:productmain}. Together with the upper bound already proved,
this establishes the claimed $\Theta(\min(n,\sqrt k/\eps))$ subclass rate.


\begin{corollary}[Exact trace recovery on the subclass is optimally $n$ calls]
If the desired error is exactly zero, a randomized algorithm that recovers the
trace of every PSD product input with probability at least $2/3$ must use at
least $n$ calls in the worst case. Thus Corollary~\ref{cor:exactn} is optimal
for exact trace recovery.
\end{corollary}
\begin{proof}
Use the stronger parallel oracle and the independent square Wishart prior in
local dimension $n$. With $t<n$ queries per factor, Lemma~\ref{lem:posterior}
leaves independent residuals with distributions $\chi^2_{(n-t)^2}$.
These have continuous densities even when $(n-t)^2=1$. The conditional product
$\prod_i(a_i+X_i)$ is positive and has no atoms. Any estimate determined by the
transcript and the algorithm's random seed therefore has conditional probability
zero of being exactly equal to the trace product. The average success is zero,
contradicting a uniform success guarantee. The argument transfers to the weaker
original product oracle.
\end{proof}

\section{Blocking and a general lower bound}
\label{sec:blocking}

\begin{proof}[Proof of the lower bound in Theorem~\ref{thm:main}]
Fix $1\le b\le q$ and let $k=\floor{q/b}$. Partition the first $bk$ modes into
$k$ blocks of size $b$, leaving $q-bk$ modes. Restrict the general PSD input to
\[
 M=W_1\otimes\cdots\otimes W_k\otimes I_{n^{q-bk}},
       \qquad W_j\in\R^{n^b\times n^b},\quad W_j\succeq0.
\]
For the lower bound, give the algorithm the individual $W_j$ matrix-vector
responses and allow arbitrary, not necessarily product, vectors of length $n^b$
in each block. One original product query is simulated by one round of this
stronger parallel oracle, using the within-block product of its original
factors. The remaining identity response is known. Relative trace estimation of
$M$ is equivalent to estimating $\prod_j\tr W_j$, because their traces differ
by the known factor $n^{q-bk}$.

Proposition~\ref{prop:lower}, with local dimension $n^b$, gives
\[
 Q(n,q,\eps)\ge\frac1{80000}
      \min\left\{n^b,\frac{\sqrt{\floor{q/b}}}{\eps}\right\}.
\]
The same constant works for each $b$; taking the maximum proves the claim.
\end{proof}

The case $b=q$ gives the ordinary-query lower scale $\min(N,1/\eps)$.
The case $b=1$ gives $\min(n,\sqrt q/\eps)$.
Intermediate blocks are useful when the local dimension is small: they prevent
the lower bound from being capped at $n$ merely because the original individual
modes are small. None of these steps assumes that an arbitrary PSD input
factors across blocks; restricting to such inputs is legitimate only on the
\emph{lower-bound} side.

\section{Scope checks, reproducibility, and remaining work}
\label{sec:scope}

\subsection{A deterministic obstruction that is not a randomized lower bound}
There is a simple deterministic $N$-call lower bound for arbitrary PSD matrices,
but it cannot be substituted for the randomized minimax lower bound.
Run a deterministic algorithm on $I_N$ using fewer than $N$ calls, and let $S$
be the span of its queried vectors. For every $c>0$, the PSD matrix
$I_N+cP_{S^\perp}$ gives exactly the same transcript: by induction all adaptive
queries stay the same and lie in $S$. Its trace is $N+c(N-\dim S)$, which can
be arbitrarily large. No common answer has fixed relative accuracy for all
these traces.

For a randomized algorithm, however, $S$ depends on its seed. Choosing a
different adversarial matrix after seeing each seed does not exhibit a single
fixed input with large failure probability. The Wishart argument in
Section~\ref{sec:lower} avoids that quantifier error by fixing an input prior
independent of the algorithm's randomness.

\subsection{What has actually been checked}
The package contains mathematical proofs, exact finite algebraic checks using
SymPy, and floating-point diagnostics using NumPy. These are different kinds
of evidence. None of the finite tests substitutes for the general arguments.

\begin{center}
\begin{tabular}{p{0.30\linewidth}p{0.60\linewidth}}
\toprule
Component & Check performed\\
\midrule
Interpolation identity & Exact symbolic vector equality for $n=2$, $q=2,\ldots,6$,
 and $n=3$, $q=2,3$.\\
Projection counterexample & Exact ranks, parity-vector norms, and rational unit-
 product projection inequalities for $q=2,\ldots,8$.\\
Parallel access & Exact rational contraction and factor-recovery checks, including
 singular PSD factors; total call count $n$.\\
Low-rank reconstruction & Exact PSD Nystr\"om identities in ambient dimension eight,
 with ranks one through four.\\
Complex second moment & $60{,}000$ Monte Carlo samples for each of three small PSD
 examples, compared with the exact partial-trace formula.\\
Adaptive Wishart residual & $7{,}000$ adaptive Krylov-query trials, with $d=10$, $t=3$,
 compared with the predicted $\chi^2_{49}$ moments.\\
\bottomrule
\end{tabular}
\end{center}

The Wishart diagnostic produced residual mean $48.98584$ and variance
$98.48736$, versus theoretical values $49$ and $98$. Its empirical correlation
with the known trace part was $-0.01174$. Agreement of these moments and a small
correlation do \emph{not} prove conditional independence; that claim rests on
the adaptive induction in Lemma~\ref{lem:posterior}.
The floating-point complex-query simulations tested through $q=6$ had relative
errors below $2.3\times10^{-15}$ on the recorded instances; no stability claim
for larger $q$ follows.

The product Hutch++ numerical example uses $n=20$, $k=2$, $r=2$, and makes
15 calls. Its particular relative error was approximately $0.1363$. This is
only a branch-execution test; $r=2$ is \emph{not} the theorem's conservative
accuracy choice for a specified small $\eps$.

\subsection{Reproduction}
From the package root, with Python, NumPy, and SymPy installed, run
\begin{verbatim}
python tests/exact_checks.py
OPENBLAS_NUM_THREADS=1 python tests/numerical_checks.py
\end{verbatim}
The scripts write machine-readable JSON results in \texttt{results/}.
The recorded numerical seed is 20260913. The reference implementations are in
\texttt{src/ra11.py}. Running \texttt{./build.sh} rebuilds this PDF with
\texttt{pdflatex}. The algorithms are written for small dense demonstrations;
large tensor dimensions, normalization underflow, and interpolation conditioning
are not engineered away.

\subsection{Final status}
The exact interpolation and parallel-access identities do not rely on numerical
experiments. The bounded-probability lower bound explicitly handles arbitrary
adaptation and removes the need to infer probability failure from RMSE alone.
The resulting subclass theorem and counterexample are concrete results of this
note, subject to independent verification.

The missing theorem is a characterization of $Q(n,q,\eps)$ for \emph{arbitrary}
PSD matrices in the full parameter range. The present upper and lower bounds
do not close that gap. Consequently this package should not be represented as
a complete resolution of RA-11, nor used by itself as grounds to mark that
problem solved.

\begin{thebibliography}{9}
\raggedright
\bibitem{ra11}
OpenProblemsInNLA, \emph{RA-11 --- Optimal trace-estimation complexity using
Kronecker matrix-vector queries}. Statement read on 13 September 2026;
repository last-check field 10 September 2026.
\url{https://github.com/ajt60gaibb/OpenProblemsInNLA/tree/main/randomized-and-low-rank-approximation/RA-11}.

\bibitem{msw}
R. A. Meyer, W. Swartworth, and D. P. Woodruff,
\emph{Understanding the Kronecker Matrix-Vector Complexity of Linear Algebra},
ICML 2025, Proceedings of Machine Learning Research 267, 43909--43933.
The theorem and conjecture numbering used here is from arXiv:2502.08029v2,
13 February 2025, in particular Theorem 7 and Conjecture 23.
\url{https://arxiv.org/abs/2502.08029v2}.
Published record: \url{https://proceedings.mlr.press/v267/meyer25a.html}.

\bibitem{ma}
R. A. Meyer and H. Avron,
\emph{Hutchinson's Estimator is Bad at Kronecker-Trace-Estimation},
SIAM Journal on Matrix Analysis and Applications 47 (2026), 353--387.
Online manuscript consulted: arXiv:2309.04952v2, 31 January 2025,
especially Sections 1.2.3, 6, and 7.
\url{https://arxiv.org/abs/2309.04952v2}.
\url{https://doi.org/10.1137/24M1720895}.

\bibitem{hutchpp}
R. A. Meyer, C. Musco, C. Musco, and D. P. Woodruff,
\emph{Hutch++: Optimal Stochastic Trace Estimation},
SOSA 2021, 142--155; arXiv:2010.09649.
\url{https://arxiv.org/abs/2010.09649}.

\bibitem{logconcave}
A. Saumard and J. A. Wellner,
\emph{Log-concavity and strong log-concavity: a review},
Statistics Surveys 8 (2014), 45--114; arXiv:1404.5886.
Used for preservation of log-concavity under convolution.
\url{https://arxiv.org/abs/1404.5886}.
\end{thebibliography}
\end{document}
